\uuid{4yNG}
\exo7id{1666}
\titre{exo7 1666}
\auteur{ortiz}
\organisation{exo7}
\datecreate{1999-04-01}
\isIndication{false}
\isCorrection{false}
\chapitre{Réduction d'endomorphisme, polynôme annulateur}
\sousChapitre{Diagonalisation}
\cpge{mpsi-mp2i.calcul-matriciel-systemes.e, pcsi.determinants.c, ptsi.determinants.a, mpsi-mp2i.calcul-matriciel-systemes.b, pcsi.calcul-matriciel-systemes-lineaires.b, ptsi.calcul-matriciel-systemes-lineaires.b, bcpst1.matrices, mp-mpi.reduction-endomorphismes.g, pc.complements-espaces-vectoriels.d, psi.complements-espaces-vectoriels.d, mp-mpi.reduction-endomorphismes.i, pc.reduction.d, psi.reduction-endomorphismes.d, bcpst2.valeurs-propres-vecteurs-propres.b}
\module{Algèbre}
\niveau{L2}
\difficulte{}

\contenu{
\texte{
Soient les matrices
$A=\left( \begin{smallmatrix}
4&0&0\\
2&3&0\\
-1&0&1
\end{smallmatrix}\right) \in \mathcal{M}_3(\Rr)$,
$D=\left( \begin{smallmatrix}
4&0&0\\
0&3&0\\
0&0&1
\end{smallmatrix}\right) \in \mathcal{M}_3(\Rr)$.
}
\begin{enumerate}
    \item \question{Soit $Y\in\mathcal{M}_3(\Rr)$ telle que $Y^2=D$.
    }
    \item \question{Montrer que $Y$ et $D$ commutent.}
    \item \question{En d\'eduire que $Y$ est diagonale puis d\'eterminer $Y$.}
\end{enumerate}
}
